% Figure: free-body diagram of a trapezoidal concrete gravity dam. Shows the
% horizontal hydrostatic thrust P_h on the upstream face at H/3, the self-weight
% W at the section centroid, the uplift U on the base, and the toe O about which
% moments are taken. No em-dashes.
\begin{figure}[htbp]
\centering
\begin{tikzpicture}[scale=0.95]
% reservoir water (left of upstream face), shaded
\fill[engblue!12] (-3.4,0) -- (0,0) -- (0,5.4) -- (-3.4,5.4) -- cycle;
\draw[engblue, thick] (-3.4,5.4) -- (0.2,5.4);
\node[engblue, font=\scriptsize, anchor=south] at (-2.6,5.4) {reservoir surface};
% trapezoidal dam cross-section: vertical upstream face, battered downstream face
% heel at (0,0), toe at (3.6,0), crest from (0,5.4) to (1.0,5.4)
\draw[engblack, very thick] (0,0) -- (0,5.4) -- (1.0,5.4) -- (3.6,0) -- cycle;
\node[engblack, font=\small] at (1.4,2.0) {dam};
% toe label O
\fill[engblack] (3.6,0) circle (1.6pt);
\node[engblack, font=\small, anchor=north] at (3.6,-0.08) {toe $O$};
\node[engblack, font=\small, anchor=north] at (0,-0.08) {heel};
% base line
\draw[enggray, thick] (-0.1,0) -- (3.7,0);
% horizontal hydrostatic thrust P_h on upstream face, acting at H/3
\draw[engamber, ultra thick, ->] (-1.7,1.8) -- (0,1.8);
\node[engamber, font=\small, anchor=east] at (-1.7,1.8) {$P_h$};
\draw[enggray, <->] (-1.95,0) -- (-1.95,1.8);
\node[enggray, font=\scriptsize, anchor=east] at (-1.95,0.9) {$H/3$};
% pressure prism on upstream face (triangular)
\draw[engamber, thin] (0,5.4) -- (0,0);
\draw[engamber, thin] (0,0) -- (-1.1,0);
\draw[engamber, thin] (-1.1,0) -- (0,5.4);
% self-weight W at section centroid, downward
\coordinate (Gd) at (1.2,1.9);
\fill[engred] (Gd) circle (1.6pt);
\draw[engred, ultra thick, ->] (Gd) -- ($(Gd)+(0,-1.6)$);
\node[engred, font=\small, anchor=west] at ($(Gd)+(0.05,-0.9)$) {$W$};
% uplift U on base, upward (distributed, shown as block of up-arrows)
\foreach \x in {0.4,1.0,1.6,2.2,2.8,3.4} {
  \draw[englime, ->] (\x,-0.9) -- (\x,0);
}
\node[englime, font=\small, anchor=north] at (1.9,-0.95) {uplift $U$};
\end{tikzpicture}
\caption[Free body of a gravity dam]{Free-body diagram of a concrete
gravity dam, drawn for the end-to-end stability check. The horizontal
hydrostatic thrust $P_h=\tfrac12\rho gH^2$ per unit width (amber) acts at
$H/3$ above the base; the self-weight $W$ (red) acts down through the
cross-section centroid; the uplift $U$ (green) acts up on the base where
pore water under head reduces the effective normal force. Moments are
taken about the downstream toe $O$: the overturning moment from $P_h$ is
resisted by $W$ acting through its lever arm to the toe, and sliding is
checked by the friction the reduced normal force $W-U$ can mobilise
against the thrust.}
\label{fig:vol03ch10:dam-fbd}
\end{figure}
